\section{Stabilizer codes}\label{sec:stabilizer_codes}

A \emph{stabilizer code} is defined by a stabilizer group $\mathcal{S}$ on $n$ physical qubits.
The \emph{codespace} is the corresponding stabilizer subspace $\mathcal{V}_{\mathcal{S}}$, which has dimension $2^k$ where $k = n - r$ and $r$ is the rank of $\SSS$.
We say $k$ is the number of \emph{logical qubits} encoded by the code.

The quotient group $\mathcal{N}(\mathcal{S})/\mathcal{S}$ is then called the \emph{logical Pauli group}.
As we saw above, this quotient group is isomorphic to $\mathcal{P}_k$, the Pauli group on $k$ qubits (hence the name!).
Cosets in this group represent logical Pauli operations on the encoded qubits: any two representatives $P$ and $Q$ from the same coset $\overline{P} = \overline{Q}$ implement the same logical operation (they differ by an element of $\mathcal{S}$, which by definition acts as the identity on the codespace).

\subsection{Implementation}\label{subsec:stabilizer_codes_implementation}

Operationally, a stabilizer code is implemented as follows.
Given $n$ qubits, we repeatedly measure a generating set $\{S_1, \ldots, S_r\}$ of the stabilizer group to project the system into the codespace and keep it there.
In the absence of any errors, each measurement $S_j$ should give the outcome $+1$ in every round, confirming that the state is still stabilized by $\mathcal{S}$.
We say each pair of measurements of $S_j$ in consecutive rounds form a \textit{detector}.

In general, a detector is a set of Pauli measurements $\{P_1, \ldots, P_d\}$ with the property that the product of their measurement outcomes is deterministic in the absence of errors.
Labelling their measurement outcomes as $m_1, \ldots, m_d$, we'll write detectors as \textit{formal products} $D = m_1 \ldots m_d$ with powers modulo $2$.
I personally think writing these as formal products is incredibly convenient, but I seem to have found myself in a bit of a minority there.
So an explainer is perhaps in order.

`Formal' just means we're viewing $m_j$ as a variable rather than a number $\pm 1$, much like when we write something like $x^2 + x + 1$.
`Product' means we are allowed to multiply these variables $m_j$ and $m_k$ together to get $m_j m_k$, which we can in turn multiply with other such expressions, and so on.
It also lets us use the variable $1$, which behaves as expected: $1 \cdot m_j = m_j = m_j \cdot 1$.
We can use power notation for multiplication: $m_j m_j \eqqcolon m_j^2$.
That we're taking `powers modulo 2' just means $m_j^n \coloneqq m_j^{n \text{ mod } 2}$, so for example $m_j^2 = m_j^0 = 1$.
The reason we enforce this is that ultimately these formal products will be \textit{evaluated} by assigning numbers $\pm 1$ to each variable $m_j$ (just like evaluating a polynomial $f(x_1, \ldots, x_d)$).
Since $(\pm 1)^2 = 1$, we enforce this at the formal level too.
We note that the set of all detectors form a group $\DDD$ under this product; if $D = m_1 \ldots m_d$ and $D' = m_1' \ldots m'_{d'}$ are detectors, then by definition in the absence of noise they evaluate to determinstic values $\pm 1$.
But then certainly the product of these two values is again deterministic, so $DD' = m_1 \ldots m_d m_1' \ldots m'_{d'}$ is another detector.
We denote the group of all detectors on a given circuit by $\DDD$.

Explainer over -- let's get back to it.
If a detector $D = m_1 \ldots m_d$ deterministically evaluates to $\pm 1$ in the absence of errors, but during a given implementation of the code evaluates to $\mp 1$ instead, then we say the detector has been \textit{violated}.
The set of all violated detectors during this implementation of the code is the \textit{syndrome}.
The syndrome is what's passed to a \textit{decoder}.
A decoder is a classical algorithm whose job is to figure out what errors are likely to have caused these detectors to be violated, then correct these errors.
We will discuss this later in \Cref{sec:decoding}.

\begin{example}\label{ex:422_stabilizer_code}
    Consider $\SSS = \langle X_1 X_2 X_3 X_4, Z_1 Z_2 Z_3 Z_4 \rangle \leq \PPP_4$, as in \Cref{ex:422_stabilizer_group_updates}.
    Viewed as a stabilizer code, this is called the $\qcode{4, 2, 2}$ code (it uses $4$ physical qubits to encode $2$ logical qubits with an \textit{algebraic distance} of $2$; more on distance later).
    It has logical Pauli group:
    \begin{equation}
        \mathcal{N(S)/S} = \langle
            \overline{i},
            \overline{\XXX_1} = \overline{X_1 X_2}, \overline{\ZZZ_1} = \overline{Z_1 Z_3},
            \overline{\XXX_2} = \overline{X_1 X_3}, \overline{\ZZZ_2} = \overline{Z_1 Z_2}
        \rangle
    \end{equation}
    To implement this stabilizer code, we can just repeatedly measure the stabilizer generating set of $X_1 X_2 X_3 X_4$ and $Z_1 Z_2 Z_3 Z_4$.
    As we saw in \Cref{ex:422_stabilizer_group_updates}, after the initial $0$th round of measurements, every measurement will be deterministic.
    That is, every such measurement will form a detector.
    In the notation of \Cref{ex:422_stabilizer_group_updates}, every $a_j a_{j+1}$ and $b_j b_{j+1}$ will be a detector, for all $j \geq 0$.
\end{example}

\begin{example}\label{ex:colour_code_stabilizer_code}
    Suppose we place a honeycomb lattice on a torus, as in Figure TODO.
    Let $V, E, F$ denote the vertices, edges and faces of this lattice respectively.
    The honeycomb lattice has the property that it's three-valent; every vertex has degree three.
    It's also three-face-colourable; every face can be assigned one of three colours such that no faces that share an edge have the same colour.
    (In fact, we can define the colour code on any planar graph with this property -- it need not be the honeycomb code specifically.)
    Edges then inherit a colouring -- an edge $e$ that borders faces $f_{\kappa}$ and $f_{\kappa'}$ of colours $\kappa \neq \kappa'$ respectively is assigned the third colour $\kappa''$.
    Throughout this thesis, we use colours cyan ($\colourful{c}$), magenta ($\colourful{m}$) and yellow ($\colourful{y}$).
    So we can partition faces $F$ into $F = F_{\colourful{c}} \cup F_{\colourful{m}} \cup F_{\colourful{y}}$, and edges $E = E_{\colourful{c}} \cup E_{\colourful{m}} \cup E_{\colourful{y}}$.

    Now, at every vertex $v \in V$, suppose we place a qubit, and for every face $f_\kappa \in F$ define weight-six \textit{face operators} $\dcccOp{f}{\kappa}{X} = \prod_{v \in f_{\kappa}} X_v$ and $\dcccOp{f}{\kappa}{Z} = \prod_{v \in f_{\kappa}} Z_v$, where $v \in f$ means $v$ is a vertex of face $f$.
    Define:
    \begin{equation}
        \SSS = \pres{\dcccOp{f}{\kappa}{X}, \dcccOp{f}{\kappa}{Z} : f_{\kappa} \in F} \leq \PPP_{|V|}
    \end{equation}
    This stabilizer group $\SSS$ defines the \textit{colour code}, a well-known static stabilizer code.
    On a torus it has logical Pauli group $\NNN(\SSS)/\SSS \cong \PPP_4$, so encodes $4$ logical qubits.
    For any edge $e = (u, v)$, define \textit{edge operators} $\dcccOp{e}{\kappa}{X} = X_u X_v$ and $\dcccOp{e}{\kappa}{Z} = Z_u Z_v$.
    We can choose representatives $\XXX_j$ and $\ZZZ_j$ for a generating set of cosets of $\NNN(\SSS)/\SSS$ to be certain `loops' of these edge operators.
    That is, we can arbitrarily pick two colours $\kappa$ and $\kappa'$, and consider `loop' operators as shown in Figure TODO for Paulis $X$ and $Z$.
    We have two Paulis $X$ and $Z$, two non-trivial loops on a torus, and two colours; each of these eight combinations gives one of $\XXX_1, \ZZZ_1, \ldots, \XXX_4, \ZZZ_4$.

    To implement this code, we again just measure the given generators of $\SSS$ -- here it's the $X$ and $Z$ face operators.
    After the initial $0$th round of measurements, every such measurement is deterministic, so forms a detector.
\end{example}
